\section{模型的唯一性及连续残差能说明什么}
\subsection{界面法向与分部积分}
界面条件同时连接两个区域，法向约定必须在模型、离散和误差计算中一致。设$n$从$\Omega_1$指向$\Omega_2$，则两个子域在界面上的外法向分别为$n$和$-n$。对于连续试验函数$v$，两侧分部积分产生的界面项为
\[
\int_\Gamma \beta_1\partial_n u_1v-\int_\Gamma\beta_2\partial_n u_2v
=-\int_\Gamma g_qv.
\]
因此，若解连续、外边界为零，经典形式对应
\begin{equation}\label{eq:weak_interface}
\sum_{i=1}^2\int_{\Omega_i}\beta_i\nabla u_i\cdot\nabla v
=\sum_{i=1}^2\int_{\Omega_i}f_iv-\int_\Gamma g_qv.
\end{equation}
这个负号来自本文的跳跃定义，而非通量条件本身具有另一种物理含义。若改用左减右的跳跃约定，右端符号随之改变。强形式配置不需要实际计算这些积分，但分部积分提供了独立检查通量符号的途径。

非零值跳跃时，解不属于通常的全域连续空间，不能直接把$u$当作齐次Dirichlet能量空间中的函数。本文使用两侧独立特征，正是为了保留这个分片性质。弱形式仍可通过适当的跳跃提升处理；不同弱离散如何惩罚值跳跃属于另一组方法选择，不等同于本文的强配置约束。

\begin{proposition}[直线界面经典解的唯一性]\label{prop:unique}
假设两组解在每侧闭区域上二阶连续，满足相同的内部、外边界、值跳跃和通量跳跃数据。扩散系数为正常数。则两组解相同。
\end{proposition}
\begin{proof}
令差为$w_i$。它满足$-\beta_i\Delta w_i=0$、外边界为零、$[w]=0$及$[\beta\partial_nw]=0$。两侧分别乘$w_i$并分部积分；外边界项为零，界面上$w_1=w_2$，两侧通量项之和也为零。故
\[
\sum_i\beta_i\int_{\Omega_i}|\nabla w_i|^2=0.
\]
由$\beta_i>0$知每侧梯度为零，因而每侧为常数。连续界面使常数相同，零外边界使常数为零。
\end{proof}
该命题只证明在给定光滑类别中的唯一性，没有单独证明任意粗糙数据下的存在性。本文通过先指定分片光滑函数再制造数据，已经给出了所用算例的存在性，因此无需用更强的一般存在理论替代实际模型核查。这个证明也说明，漏掉值条件或通量条件会改变问题，而不能仅视为少加一些训练行。

\subsection{带精确边界和值条件的条件性残差界}
为了区分连续残差与有限训练损失，考虑一个比实际计算更理想的情况。设真解$u$与近似$v$均分片光滑，并且$v$精确满足外边界和值跳跃。记$e=v-u$，连续内部残差$r_i=-\beta_i\Delta v_i-f_i$，连续通量残差$s=[\beta\partial_nv]-g_q$。这里连续范数使用区域面积及界面长度，而不是训练点的均方。

由于$e$在界面连续且在外边界为零，对固定$y$可沿$x$积分：$e(x,y)=\int_0^x e_x(t,y)\,dt$，积分经过界面时由连续性连接两段。Cauchy--Schwarz不等式给出
\[
\norm e_{L^2(\Omega)}\le\norm{\nabla e}_{L^2(\Omega)},\qquad
\norm{e|_\Gamma}_{L^2(\Gamma)}\le 2^{-1/2}\norm{\nabla e}_{L^2(\Omega)}.
\]
这些常数是为当前正方形给出的简便上界，不主张为最优常数。
\begin{proposition}[特定假设下的连续残差控制]\label{prop:energy}
在上述精确边界和值跳跃条件下，令$\beta_{\min}=\min(\beta_1,\beta_2)$，则
\begin{equation}
\norm{\nabla e}_{L^2(\Omega)}
\le\frac{\norm r_{L^2(\Omega)}+2^{-1/2}\norm s_{L^2(\Gamma)}}{\beta_{\min}}.
\end{equation}
其中$\norm r^2=\sum_i\norm{r_i}_{L^2(\Omega_i)}^2$。
\end{proposition}
\begin{proof}
对误差方程分部积分，得到
\[
\sum_i\int_{\Omega_i}\beta_i|\nabla e_i|^2
=\sum_i\int_{\Omega_i}r_ie_i-\int_\Gamma se.
\]
左端不小于$\beta_{\min}\norm{\nabla e}^2$，右端绝对值不大于$(\norm r+2^{-1/2}\norm s)\norm{\nabla e}$。若梯度范数为零则结论显然，否则约去该范数。
\end{proof}
这只是能量法在受限情形下的自含推导，不是本文自适应配置点算法的一般后验估计。实际网络仅在有限点弱满足外边界和值跳跃，因而不能直接把表中的离散残差代入本命题。若还存在值跳跃误差，需要构造提升并控制其导数；单独知道界面上几个点的值误差不足以完成这一步。

例如，值跳跃误差为光滑函数$a(y)$且$a(0)=a(1)=0$时，可取左侧提升为零、右侧提升$\ell_2=2(1-x)a(y)$。它在界面产生值跳跃$a$并保持外边界为零，但同时产生通量跳跃$-2\beta_2a$以及含$a''$的内部项。由此可见，处理值误差还涉及沿界面的正则性，不能只把一个值RMS与通量RMS相加就称为完整能量估计。本文联合指标是一种可计算的启发式量，其数学作用应限定在这个层面。

\section{随机特征空间、可辨识性与观测范数}
\subsection{随机与线性的两个含义}
固定隐藏参数之后，函数集合
\[
V_N=\operatorname{span}\{\phi_1,\ldots,\phi_N\}
\]
是有限维线性空间。随机性来自生成这个空间的参数，而不是输出系数求解本身。对于给定空间和给定点集，线性PDE对应确定的矩阵问题；不同随机种子改变空间以及点位置，实验离散程度因此包含不止一种随机来源。共享特征能消除方法间的一部分差异，但不能让不同种子变成完全相同的近似问题。

特征个数不一定等于空间维数。若两个特征完全相同，其系数可以相互抵消；即使数学上线性独立，也可能在有限精度或特定配置点上近乎相关。有效秩衡量的是已加权观测矩阵可以分辨的方向，它既不同于形式上的特征数，也不同于连续函数空间的精确维数。这解释了为什么本文同时报告总列数和截断后的数据秩。

采用分片空间时，$V_N=V_{N_1}^{(1)}\oplus V_{N_2}^{(2)}$中的函数在界面可以不连续。界面行把两个空间连接起来，而不是预先要求所有基函数连续。对非零值跳跃，这种表示非常直接；对零跳跃，连续性由方程约束实现。若删除界面行，两个子域求解将失去所需的耦合，内部残差小也不能单独恢复原模型。

\subsection{有限维观测控制的必要条件}
设$L$把$V_N$中的函数映射为所有加权内部、边界和界面观测。只有$L$在所考虑空间上为单射时，$\norm{Lv}$才是一个范数；若存在非零$v$使所有观测为零，它只能是半范数。有限点配置的稳定性因此首先涉及可辨识性，而不是简单地把点数加到大于特征数。

在系数表示下，若$B$满列秩，对任何预测矩阵$P$有
\begin{equation}\label{eq:prediction_bound}
\norm{Pc}_2\le\norm{PB^\dagger}_2\norm{Bc}_2
\le\frac{\norm P_2}{\sigma_{\min}(B)}\norm{Bc}_2.
\end{equation}
第一步使用$c=B^\dagger Bc$。该式是固定空间、固定观测和固定验证点的确定性结论；它提示预测稳定性常数$\norm{PB^\dagger}$比仅看$\kappa(B)$更接近具体函数预测任务，但本文尚未对全部实验计算这个常数。

若矩阵截断后，只能对保留右奇异空间上的系数使用对应界。丢弃空间中的成分不受这个训练观测界的直接控制，第二章的误差分解已将其单列。相对阈值改变时，保留空间可能变化，不能假设所有阈值下比较的是同一个近似对象。这也解释了为什么截断与配置点的交互必须结合奇异值实际检查。

\subsection{有限点小残差不等于连续小残差}
一个初等反例能说明有限观测的逻辑边界。给定区间内有限点$x_1,\ldots,x_m$，多项式$\psi(x)=\prod_{j=1}^m(x-x_j)$在这些点全部为零，但在点集之外通常不为零；任意倍数$C\psi$仍具有零经验观测，连续范数却随$|C|$增大。这个反例不是说本文tanh空间一定包含该多项式，而是说明若缺少空间限制及稳定观测假设，不能从有限点数值推出一般连续结论。

对于固定有限维特征空间，充分覆盖和矩阵可辨识性可以排除某些这种自由度，但常数仍可能很大。对于不断改变的特征或自适应点集，则需进一步控制空间和采样的共同变化。本文没有证明随机点集以何种概率满足统一范数等价，因此所有采样结论以独立离散验证和重复种子为证据，不写成未建立的概率收敛定理。
